\documentclass[11pt]{article}
\usepackage[margin=0.95in]{geometry}
\usepackage{amsmath,amssymb,amsthm,booktabs}
\usepackage[hidelinks]{hyperref}
\setlength{\emergencystretch}{2em}
\newtheorem{theorem}{Theorem}
\newtheorem{lemma}[theorem]{Lemma}
\newcommand{\mr}{\operatorname{mr}}
\newcommand{\rank}{\operatorname{rank}}
\title{The original sixteen belong to a known graph-cover MRC family}
\author{Source-scope audit and classical rigidity consequence}
\date{17 November 2026}
\begin{document}
\maketitle

The original sixteen working equalities remain valid. A further classical
rigidity step now identifies the actual double branched covers with the
constructed graph-manifold models, up to orientation. Consequently all
sixteen knots belong to the family already covered by
Boileau--Jang--Weidmann's 2018 meridional-rank theorem. These examples
must not be described as cases outside known MRC families. The explicit
group ranks and bridge-number certificates may still be useful; their
named global priority has not been established. Full MRC remains open.

\begin{theorem}\label{thm:family}
For every original target $K$ with index in
\[
129,135,136,137,194,196,197,198,233,234,239,242,243,244,253,265
\]
in the $13a$ table, the actual manifold $\Sigma_2(K)$ is a closed prime
graph manifold. The preceding working proofs give
\[
 \rank\pi_1(\Sigma_2(K))=3,\qquad b(K)=\mr(K)=4.
\]
Hence the MRC equality lies within the scope of the published
Boileau--Jang--Weidmann graph-cover theorem.
\end{theorem}

\section{The retained input and the additional rigidity step}

For each named source, the preceding exact core and peripheral proofs
construct a closed orientable irreducible graph manifold $M_K$ and an
isomorphism
\[
 \pi_1(\Sigma_2(K))\cong\pi_1(M_K).
\]
The constructed manifolds have three or four canonical Seifert pieces,
with primitive injective torus maps and nonparallel unique fibre slopes.
Each contains an injected group $\langle a,t\mid t^2=a^3\rangle$, whose
central element is infinite. Thus $M_K$ is prime, has infinite fundamental
group, and is neither a lens space nor $S^1\times S^2$.
The previous rank proofs needed only these group models. We now use a
separate classical input to obtain actual-cover topology.

Aschenbrenner--Friedl--Wilton's \emph{3-Manifold Groups}, final author
manuscript, Theorems~2.1.2--2.1.3, classify closed orientable
3-manifolds with isomorphic fundamental groups. The remaining freedom
is orientation of prime factors and lens-space ambiguity. Apply the
arbitrary-closed statement to $\Sigma_2(K)$ and $M_K$. The latter is
prime with infinite group, so it has no nontrivial lens summand and only
one nontrivial prime factor. The corresponding source decomposition
therefore gives a homeomorphism to $M_K$, possibly reversing orientation.
This also proves source-cover primeness; it is not assumed from a knot
table label. Graph-manifold membership is unaffected by orientation.

The imported statement is read on final author manuscript printed
pages~30--31 (PDF pages~34--35), with the preceding prime-decomposition
discussion. Its earlier author-v3 counterparts, Theorems~2.2--2.3, were
compared at the same scope. Publisher metadata identifies the 2015 EMS
book, \url{https://doi.org/10.4171/154}. Original constituent
geometrization/Mostow/Waldhausen/Scott proofs and the full published
version are not independently audited here.

\section{The existing MRC theorem and what it already covers}

The published Boileau--Jang--Weidmann paper, \emph{Meridional rank and
bridge number for a class of links}, Pacific Journal of Mathematics~292
(2018), \url{https://doi.org/10.2140/pjm.2018.292.61}, defines a graph
manifold to be prime. Theorem~1.2 and Corollary~1.3 on printed page~62
therefore apply to the actual covers just identified: meridional rank
three forces bridge three in this cover class, and exact bridge four
forces meridional rank four. The theorem treats arborescent and
non-arborescent links, so proving an arborescent tree description is not
necessary for this family-membership conclusion. The full proof remains
an explicit classical import.

Together with the standard rank-one and rank-two meridional equivalences
credited in that introduction, these statements give a useful scope
consequence: any knot in the graph-cover class with a certified upper
$b(K)\le4$ already satisfies MRC. Indeed $\mr(K)\le4$; for meridional
rank one, two or three the corresponding bridge number is the same,
while meridional rank four and the upper force bridge four. Thus the
equality itself can lie in an existing class even before deciding
whether its exact value is three or four.

Our cover-rank-three calculation does decide the value: the
even-meridian bound gives $\mr(K)\ge4$, and each retained source has an
individual bridge upper four. The family audit preserves that stronger
explicit input. It does not establish whether the particular named
bridge/rank computations appeared earlier elsewhere.

\section{Exact case bindings and verification limits}

The retained cases bind as follows. Identical homology orders in
different rows are not used to identify groups or branch sets.
\begin{center}
\begin{tabular}{lll}
\toprule
Indices & Retained model route & JSJ pieces\\
\midrule
$129,136,194,197$ & root $129$ and certified mutation aliases & $4$\\
$135,137,196,198$ & root $135$ and certified mutation aliases & $4$\\
$233,239,253$ & three direct peripheral models & $3$\\
$234,242,243,244,265$ & five direct peripheral models & $3$\\
\bottomrule
\end{tabular}
\end{center}
The first two rows use only the six previously checked mutation/common-cover
certificates, never mutation invariance of bridge or meridional rank.
The other eight cases use their own exact source groups. All sixteen
individual upper-four traces are rebound to their literal retained PDs.
The companion scope receipt checks source hashes, ledger dependencies,
the exact sixteen-name set and these finite bindings. It does not
algorithmically prove manifold rigidity or the published family theorem.

The lens-space exception matters: $L(5,1)$ and $L(5,2)$ share group
$C_5$ yet are not homeomorphic. Also the same rigidity shortcut is not
valid for arbitrary manifolds with boundary without peripheral data.
Here both the actual covers and models are closed, and the infinite
prime models remove the relevant exceptions. No homeomorphism respecting
the branched-cover involution is needed: the known theorem asks only for
the actual cover's graph-manifold type. No numerical triangulation,
rank-equals-Heegaard-genus assumption, or branch-set classification is used.

The three earlier exact case proofs and their data are retained. Source
catalogue46 adds one scoped rigidity reference and a published snapshot
of the existing graph-cover reference. Refreshing affected hashes used
only the already frozen small recipes and their replays; no293-knot or
base/module campaign was repeated.

\section{The next research boundary}

The original sixteen are now established members of an existing MRC
family, rather than merely candidates of unknown family status.
The separate exact topology audit in \\path{SEED_COVER_TOPOLOGY.tex}
proves that the four earlier simple-three-cover seeds lie outside this
graph-cover class, by incompressible tori cutting off certified hyperbolic
pieces. The arbitrary-cycle bound in
\\path{GENERAL_BRANCHED_COVER_RANK_BOUNDS.tex} also supplies twenty-four
further sharp-four equalities, all outside the same graph-cover class.
Thus forty-four of the retained293 names have current certificates.
Global named priority, other-family and positive periodic-parameter
comparisons remain open, as do the other249 panel names and full MRC.

For progress beyond the BJW18 family, a new target must have separately
justified cover topology outside this class, or another route whose
hypotheses are not confined to graph-manifold double covers. A hyperbolic
knot exterior alone does not exclude a graph-manifold double cover.
Likewise non-arborescence alone is insufficient, since the known theorem
also treats that branch. A quotient of the double-cover group with ordinary
rank at least three can still give a lower four without identifying the actual cover
with a graph manifold; constructing such a quotient on a justified
outside-family target is a distinct prospective task, not a result here.

The unpublished affine argument's bridges to signed lifts and cover
homology remain proved at their stated scopes, together with its retained
universality obstruction. This family audit does not expand that affine
criterion. Full MRC, exhaustive literature, formal and external specialist
verification, and exact named computational priority remain open.

\end{document}
